\begin{textsample}{POS Dim 3 – ms – Score 8.00 – ms/ms\_inf\_e\_ef\_24.txt}  \label{ex:f3_pos_238}
7.1 Infância e Criança Cresci \textbf{brincando} no \textbf{chão}, entre formigas.
De uma infância livre e sem comparamentos.
Eu tinha mais comunhão com as coisas do que comparação.
Porque a gente fala a partir do ser \textbf{criança}, a gente faz comunhão : de um orvalho e sua aranha, de uma tarde e suas garças, de um pássaro e sua árvore. ( Manoel de Barros ) Compreende -se que a infância é um conceito construído socialmente, fruto do desenvolvimento histórico da humanidade, portanto uma categoria histórica e social.
Pode -se entendê -la, ainda, “ como período da história de cada um, que se estende, na nossa sociedade, do nascimento até aproximadamente doze anos de idade ” ( KRAMER, 2003, p. 32 ).
Portanto, a necessidade de compreender a infância e sua especificidade exige caracterizar a \textbf{criança} concreta e historicamente, e assumi -la como sujeito e cidadã de direitos, que se constitui na sociedade da qual faz parte.
De acordo com Resolução n. 5/2009/CEB/CNE, art. 4º : [... ]"a \textbf{criança}, centro do planejamento curricular, é sujeito histórico e de direitos que, nas interações, relações e práticas cotidianas que vivencia, constrói sua identidade pessoal e coletiva, \textbf{brinca}, imagina, fantasia, deseja, aprende, observa, experimenta, narra, questiona e constrói sentidos sobre a natureza e a sociedade, produzindo cultura".
Assim, na Educação \textbf{Infantil} a \textbf{criança} deve vivenciar cotidianamente essas ações e experiências para que possa aprender, se desenvolver e produzir cultura.
Estudos de Dahlberg, Moss e Pence ( 2003 ) explicitam diferentes entendimentos do que seja e do que deva ser \textbf{criança} na sociedade multifacetada atual.
Com as concepções atribuídas à infância no âmbito da educação institucionalizada planejam -se as experiências, organizam -se os espaços e os tempos que influenciam diretamente o processo de aprendizagem e desenvolvimento e, consequentemente, a formação cultural das \textbf{crianças}.
Mesmo que não se tenha consciência dessas ideias e concepções, elas direcionam as práticas, ações e intenções.
Dahlberg, Moss e Pence concebem e apresentam uma concepção de \textbf{criança} como co-construtora de conhecimento, identidade e cultura, um sujeito que tem vez e voz : > A infância é uma construção social, elaborada para e pelas \textbf{crianças}, em um conjunto ativamente negociado de relações sociais ; A infância como construção social é sempre contextualizada em relação ao tempo, ao local e à cultura, variando segundo a classe, o gênero e outras condições socioeconômicas.
Por isso, não há uma infância natural nem universal, e nem uma \textbf{criança} natural ou universal, mas muitas infâncias e \textbf{crianças} ( 2003, p. 71 ).
Portanto, a ideia da infância única, abstrata e desvinculada da dinâmica da sociedade não pode ser sustentada, pois as \textbf{crianças} são seres com histórias e experiências diferentes, tendo, por conseguinte, mecanismos, necessidades e condições diferenciadas para fazerem parte do meio social onde vivem.
Desse modo, no documento referência para a construção do currículo, busca -se uma proposta que considere as especificidades da infância, os modos de ser das \textbf{crianças} e suas infâncias, os saberes constituídos por elas no cotidiano a fim de articulá -los com o patrimônio histórico, social e cultural construído pela humanidade.
Nessa perspectiva, concebe -se a \textbf{criança} como um sujeito de direitos, defendendo a atenção aos seus interesses e necessidades desde o nascimento.
Considerá -la como sujeito de direitos implica, dentre outros, garantir no currículo da Educação \textbf{Infantil} a presença de experiências que lhe oportunize conviver, \textbf{brincar}, participar, explorar, expressar e conhecer -se.
Considerando essas formas pelas quais \textbf{meninos} e \textbf{meninas} de zero a cinco anos aprendem, é necessário, ainda, pensar e estruturar o trabalho educativo a partir de uma concepção de \textbf{criança} que pensa, age, deseja, narra, imagina, questiona e elabora ideias sobre o mundo e, assim, criar oportunidades para ela se manifestar por meio de diferentes linguagens.
Também é necessário constituir acessibilidade de espaços, materiais, objetos, propostas, \textbf{brinquedos} e instruções para as \textbf{crianças} com deficiência, transtornos globais de desenvolvimento e altas habilidades/superdotação, conforme preconiza as Diretrizes Curriculares Nacionais para a Educação \textbf{Infantil} ( 2009 ).
A \textbf{criança} se desenvolve amplamente a partir da sua própria atividade mediante as relações estabelecidas com os outros seres humanos, as parcerias que se formam nas relações com os \textbf{adultos} e com outras \textbf{crianças} e as condições adequadas de vida e de educação.
Nas primeiras décadas do século XX, Vigotski ( 1998 ) já havia apontado que a \textbf{criança}, desde bem \textbf{pequena}, é capaz de estabelecer relações com as outras e com as coisas, em um processo em que percebe os significados e se apropria dos mesmos, atribuindo -lhes sentido.
Dessa forma, a Educação \textbf{Infantil} como primeira etapa da Educação Básica contribui qualitativamente para o desenvolvimento dos aspectos intelectuais, emocionais, físicos, sociais e culturais da \textbf{criança}.
Com base nos aprendizados que realiza, amplia suas possibilidades de atividade prática e intelectual e apropria -se, nesse processo, de novos conhecimentos.
Isso, significa compreendê -la como fruto das relações sociais que geram as diversas formas de vê -la e produzir a consciência da particularidade \textbf{infantil}.

% matched lemmas: adulto, brincar, brinquedo, chão, criança, infantil, menino, pequeno
\end{textsample}
